\section{Describe what changed between frames}
A lossless animation can store every frame as a separate image. That gives a useful
starting point: every frame is independent, and decoding it requires no earlier
pixels. It also repeats a great deal of data when only a small region moves. This
chapter changes the representation so an unchanged pixel can occupy part of a run
count rather than another RGB triple.

Build a small codec for 64-by-64 RGB images. Periodic full frames establish independent
starting points. Between them, XOR deltas describe changes to the preceding decoded
frame. A skip/change command stream groups unchanged and changed pixels. After the
scalar version is correct, use MSA to compute several pixel deltas and count changed
pixels together. Measure encoded bytes and simulated work separately.

The teaching format defined here is called \texttt{MLV1}. Its viewer implements the
byte layout in this chapter. The fixed geometry keeps the project focused on frame
dependencies, bounded command interpretation, seeking and publication. The following
chapter extends PNG with APNG's standardized animation rules. Keep this chapter's
own format rules explicit rather than treating a convenient encoder output as the
only legal representation.
